\documentclass[10pt,a4paper]{amsart}
\usepackage[margin=19mm]{geometry}
\usepackage{amsmath,amssymb,mathtools}
\usepackage[scr=rsfs,cal=euler]{mathalfa}
\usepackage{xfrac}
\usepackage[unicode,hidelinks,bookmarks=false]{hyperref}
\newcommand{\ie}{i.e.\ }
\newcommand{\cf}{cf.\ }
\newcommand{\resp}{respectively\ }
\newcommand{\Subsection}{\subsection}
\providecommand{\nobreakdash}{\nobreak\hbox{-}}
\setlength{\emergencystretch}{2em}
\multlinegap0pt
\usepackage{bm}
\usepackage{bbm}
\usepackage{dsfont}
\usepackage{xspace}
\usepackage{cancel}
\usepackage{mathtools}
\usepackage{amsthm}
\makeatletter
\@ifundefined{theorem}{\newtheorem{theorem}{Theorem}}{}
\@ifundefined{lemma}{\newtheorem{lemma}[theorem]{Lemma}}{}
\makeatother
\title[How Out-of-Distribution Detection Learning Theory Enhances T]{How Out-of-Distribution Detection Learning Theory Enhances Transformer: Learnability and Reliability}
\author{Yijin Zhou, Yutang Ge, Wenyuan Xie, Linqian Zeng, Xiaowen Dong, Yuguang Wang}
\date{Source: arXiv:2406.12915v6 (CC BY 4.0)}
\begin{document}
\maketitle

\noindent\emph{Source and licence.} How Out-of-Distribution Detection Learning Theory Enhances Transformer: Learnability and Reliability. Version arXiv:2406.12915v6 (CC BY 4.0), distributed under the Creative Commons Attribution 4.0 International licence, as recorded in the arXiv licence metadata for this exact version and shown on the abstract page https://arxiv.org/abs/2406.12915v6. Full rights notice: licenses/arxiv-2406.12915-CC-BY-4.0.txt.\par\smallskip

\noindent\textit{Editorial scope.} Counted unit: the Lemma whose source label is \texttt{Lemma 5} in the article (source0.tex lines 1007--1010, source label \texttt{Lemma 5}), delivered together with the article's own complete original proof of it (source0.tex lines 1013--1070, one \texttt{proof} environment of 58 lines). No mathematics is written, repaired or re-derived: statement and proof are the article's own text line for line. Required context retained uncounted: Jackson lemma1 (lines 944--948); eq: 5.4 (lines 951--955); eq: 5.5 (lines 969--972); Lemma 4 (lines 992--995). Every definition, display and label of those lines is retained unchanged. Omitted from this package and not redistributed: 1--943, 949--950, 956--968, 973--991, 996--1006, 1011--1012, 1071--1758. Nothing inside a retained excerpt is omitted or altered: every retained body is delivered exactly as the article's lines are written, so no omission receipt of any kind accompanies this package. Citations: the retained excerpts carry no citation command at all, so no bibliography entry of the article is transcribed and no reference is redistributed. Reference adaptation: none was needed and none was made; every reference inside the delivered unit resolves inside it, so no unresolved reference or citation is produced. Printed slips kept literal: no printed word, symbol, hypothesis or number of the counted statement or of its proof was altered. Layout and class: the article's own class is \texttt{article} with packages including microtype, graphicx, subcaption, booktabs, hyperref, icml2026, inputenc, fontenc, url, amsfonts, nicefrac, xcolor, amsmath, amssymb; the wrapper is the lane's pinned \texttt{amsart} editorial wrapper at 10pt on A4, which restates the theorem-like environments the retained text needs and adds the article's own macro definitions that the retained text needs (none), with extra wrapper packages (bm, bbm, dsfont, xspace, cancel, mathtools). The delivered numbering is the unit's own. Assets: no retained span contains a figure environment, a graphics-inclusion command, a PDF-page inclusion, a table or a TikZ picture; no image file is copied into this package, the package declares no asset dependency and no image file is redistributed with it. Licence: version arXiv:2406.12915v6 of this article is distributed under the Creative Commons Attribution 4.0 International licence, as recorded in the arXiv licence metadata for this exact version and shown on the abstract page https://arxiv.org/abs/2406.12915v6. A full rights notice is delivered at licenses/arxiv-2406.12915-CC-BY-4.0.txt. Characters: every retained excerpt is pure ASCII, so the delivered engine, XeLaTeX with its default Latin Modern text font, reports no missing character.
\begin{lemma}
    For any $\mathbf{h}\in \widetilde{\mathcal{C}}^{(\alpha, \beta)}$, and any compact set $C\in \mathbb{R}^d$, there exists a two layer transformer $\hat{\mathbf{H}} \in \mathcal{H}_{\rm Trans}^{(m,2)}$ and a linear read out $\mathbf{c}:\mathbb{R}^{\hat{d}\times \tau}\xrightarrow{}\mathbb{R}^{1\times \tau}$ \textit{s}.\textit{t}. the inequality \eqref{eq: 5.5} is established,
    where $m=(2m_h+1,1,m_h,2\tau \hat{d}_0+1,r)$.
    \label{Jackson lemma1}
\end{lemma}

    \begin{equation}        \int_{\mathcal{I}}\sum_{t=1}^{\tau}|\mathbf{c}\circ\mathbf{H}_t(\mathbf{x})-\mathbf{h}_t(\mathbf{x})|d\mathbf{x}
        \leq \tau^2 C_0(\mathbf{h} )
        \left(\frac{C_1^{(\alpha)}(\mathbf{h} )}{m_h^{2\alpha-1}}+\frac{C_2^{(\beta )}(\mathbf{h})}{m_{\rm FF}^{\beta}}(m_h)^{\beta}\right),
        \label{eq: 5.4}
    \end{equation}

    \begin{equation}   
    P(\|\mathbf{c}\circ\mathbf{H}(\mathbf{x})-\mathbf{h}(\mathbf{x})\|_2>|\mathcal{I}|(\text{RHS~in~Eq.}~\eqref{eq: 5.4}+\lambda_0))\leq\frac{\text{RHS~in~Eq.}~\eqref{eq: 5.4}}{\lambda_0|\mathcal{I}|}
    \label{eq: 5.5}
    \end{equation}

\begin{lemma}
    For any introduced $\mathbf{g}_i$ mentioned above, there exists $\hat{\mathbf{g}}_i$ satisfies $\hat{\mathbf{g}}_i \in \widetilde{\mathcal{C}}^{(\alpha, \beta)}$ and $\|\hat{\mathbf{g}}_i-\mathbf{g}_i\|_2<\epsilon$.
    \label{Lemma 4}
\end{lemma}

\begin{lemma}
    Assume the labels are independent for classification. Let $|\mathcal{X}|=n<+\infty$, $\tau>K+1$ and $\sigma$ be the {\rm Relu} function. Given any finite $\delta$ hypothesis functions $h_1,\cdots,h_{\delta}\in \{\mathcal{X}\xrightarrow{}\{1,\cdots,K+1\}\}$, then for any $m_h,r>0$, $m=(2m_h+1,1,m_h,2\tau \hat{d}_0+1,r)$, $P(h_1,\cdots,h_{\delta}\in \mathcal{H}^{(m,2)})\geq (1-\frac{m\text{RHS~in~Eq.}~\eqref{eq: 5.4}}{|\mathcal{I}|\lambda_0})^{(K+1)\delta}$ for any $\eta>1$.
    \label{Lemma 5}
\end{lemma}

\begin{proof}
    Since $\mathcal{X}$ is a compact set, then Lemma \ref{Lemma 4} implies that there exists $\hat{\mathbf{g}}_i \in \widetilde{\mathcal{C}}^{(\alpha, \beta)}$ \textit{s}.\textit{t}.
    \begin{equation}
        \|\mathbf{g}_i - \hat{\mathbf{g}_i}\|_2 < \epsilon / \|W_4\|_2.
    \end{equation}
    Denote $r_i = W_4\mathbf{g}_i^{\top}+b_4$ and $\hat{r}_i = W_4\hat{\mathbf{g}}_i^{\top}+b_4$,
    So we get
    \begin{equation}
       \|r_i -\hat{r}_i \|_2=\|W_4(\mathbf{g}_i-\hat{\mathbf{g}}_i)^{\top}\|_2\leq \epsilon.
    \end{equation}
    
    Then by Lemma \ref{Jackson lemma1}, there exists $\hat{\mathbf{H}} \in \mathcal{H}_{\rm Trans}^{(m,2)}$ and a linear read out $\mathbf{c}$ \textit{s}.\textit{t}. 
   \begin{equation}   
    P(\|\mathbf{c}\circ\mathbf{H}(\mathbf{x})-\mathbf{h}(\mathbf{x})\|_2\leq|\mathcal{I}|(\text{RHS~in~Eq.}~\eqref{eq: 5.4}+\lambda_0)\geq 1 - \frac{\text{RHS~in~Eq.}~\eqref{eq: 5.4}}{\lambda_0|\mathcal{I}|}.
    \label{eq: 5.6}
    \end{equation}
    Thus we get if $h_i(\mathbf{x})=k$, which is equal to $\mathbf{g}_i(\mathbf{x})=\mathbf{y}_k$ or $r_i(\mathbf{x})=\mathbf{z}_k$:
    
    Firstly, denote $\mathbf{f} = W_4\mathbf{c} \circ \mathbf{H}^{\top}+b_4$, and let $\mathbf{h}=\hat{\mathbf{g}_i}$, then
    \begin{equation}
        P(\|\mathbf{f}(\mathbf{x})-\hat{r}_i (\mathbf{x})\|_2\leq\|W_4\|_2|\mathcal{I}|(\text{RHS~in~Eq.}~\eqref{eq: 5.4}+\lambda_0))\geq 1 - \frac{\text{RHS~in~Eq.}~\eqref{eq: 5.4}}{\lambda_0|\mathcal{I}|}.
    \end{equation}
    So we obtain that    
    \begin{equation}
    \begin{aligned}
        &P(|\mathbf{f}_k - N|\leq\|W_4\|_2|\mathcal{I}|(\text{RHS~in~Eq.}~\eqref{eq: 5.4}+\lambda_0)) \\
        &\geq P(|\mathbf{f}_k-\hat{r}_{i,k}|+|\hat{r}_{i,k}-r_{i,k}|\leq\|W_4\|_2|\mathcal{I}|(\text{RHS~in~Eq.}~\eqref{eq: 5.4}+\lambda_0)) \\
        &\geq P(\|\mathbf{f}-\hat{r}_i\|_2+\|\hat{r}_i-r_i\|_2\leq\|W_4\|_2|\mathcal{I}|(\text{RHS~in~Eq.}~\eqref{eq: 5.4}+\lambda_0)) \\
        &\geq P(\|\mathbf{f}-\hat{r}_i\|_2+\epsilon \leq\|W_4\|_2|\mathcal{I}|(\text{RHS~in~Eq.}~\eqref{eq: 5.4}+\lambda_0)) \\
        &=P\bigr(\|\mathbf{f}-\hat{r}_i\|_2 \leq\|W_4\|_2|\mathcal{I}|\bigl(\text{RHS~in~Eq.}~\eqref{eq: 5.4}+(\lambda_0-\frac{\epsilon}{|\mathcal{I}|})\bigr)\bigr) \\
        &\geq 1-\frac{\text{RHS~in~Eq.}~\eqref{eq: 5.4}}{|\mathcal{I}|(\lambda_0-\frac{\epsilon}{|\mathcal{I}|})}\\
        &= 1-\frac{\text{RHS~in~Eq.}~\eqref{eq: 5.4}}{|\mathcal{I}|\lambda_0-\epsilon}.
    \end{aligned}
    \end{equation}
    Similarly, for any $j \neq k$, we can also obtain that
    \begin{equation}
        P(|\mathbf{f}_k|\leq\|W_4\|_2|\mathcal{I}|(\text{RHS~in~Eq.}~\eqref{eq: 5.4}+\lambda_0)\geq 1-\frac{\text{RHS~in~Eq.}~\eqref{eq: 5.4}}{|\mathcal{I}|\lambda_0-\epsilon}.
    \end{equation}
    
    Therefore, $P(\arg\max_{k\in \mathcal{Y}}\mathbf{f}_k(\mathbf{x}) = h_i(\mathbf{x}))\geq (1-\frac{\eta \text{RHS~in~Eq.}~\eqref{eq: 5.4}}{|\mathcal{I}|\lambda_0})^{K+1}$ for any $\mathbf{x}$, 
    %
    if 
    \begin{equation}
        N>2\|W_4\|_2|\mathcal{I}|(\text{RHS~in~Eq.}~\eqref{eq: 5.4}+\lambda_0)
    \end{equation}
    for any $\eta>1$, \textit{i}.\textit{e}.
    
    \begin{equation}
        P(h_1,\cdots,h_{\delta}\in \mathcal{H}^{(m,2)})\geq \left(1-\frac{\eta \text{RHS~in~Eq.}~\eqref{eq: 5.4}}{|\mathcal{I}|\lambda_0}\right)^{(K+1)\delta},
    \end{equation} 
    if
   \begin{equation}
        N>2\|W_4\|_2|\mathcal{I}|(\text{RHS~in~Eq.}~\eqref{eq: 5.4}+\lambda_0)
    \end{equation}
    for any $\eta>1$.
    Since $N$ is arbitrary, we can find such $N$.
   % \hfill\qedsymbol
\end{proof}

\end{document}
